\section*{ВВЕДЕНИЕ}
\addcontentsline{toc}{section}{ВВЕДЕНИЕ}

Рост популярности здорового образа жизни и самостоятельных занятий физической культурой сопровождается активным развитием цифровых сервисов: тренировочных дневников, трекеров активности, приложений учёта питания и социальных платформ для обмена результатами. Для пользователя, который планирует нагрузку без постоянного присутствия тренера, критичны не только справочные материалы об упражнениях, но и возможность формировать индивидуальный план, фиксировать выполнение, анализировать динамику и сохранять данные при работе с разных устройств. Эти требования естественно приводят к архитектуре <<клиент -- сервер -- база данных>>, в которой серверная часть обеспечивает аутентификацию, целостность персональной информации и единый контур обмена данными между клиентами, а клиентская часть отвечает за представление информации, ввод данных пользователем и выполнение сценариев работы в браузере.

Коммерческие решения на рынке, как правило, закрывают лишь отдельные аспекты задачи: одни продукты ориентированы на учёт силовых тренировок, другие -- на калорийность рациона, третьи -- на социальную мотивацию или экосистему носимых устройств. Объединение планирования тренировок, питания, профиля пользователя и социальной активности в одной системе при сохранении веб-доступа через браузер остаётся востребованной учебной и прикладной постановкой. В рамках настоящей выпускной квалификационной работы разрабатывается программно-информационная система <<Fitness Life>> -- бизнес-проект стартап-команды, в котором инженерный вклад автора сосредоточен на клиентской подсистеме: проектировании одностраничного интерфейса, организации модулей браузерного приложения, визуализации прогресса и взаимодействия с REST API серверной части.

Система <<Fitness Life>> реализуется как клиент-серверное веб-приложение. Серверная подсистема (разрабатывается напарником по проекту) построена на стеке Node.js и Express и взаимодействует с реляционной СУБД PostgreSQL: она принимает HTTP-запросы, проверяет учётные данные, выдаёт и проверяет JWT, хранит профили пользователей, синхронизируемое пользовательское хранилище ключ -- значение и данные социальной ленты. Клиентская подсистема представляет собой одностраничное приложение (SPA), загружаемое в браузере: HTML-разметка задаёт каркас экранов, CSS -- визуальное оформление и адаптивность, JavaScript -- интерактивную логику. После входа пользователь работает с единой оболочкой приложения и переключает разделы (тренировки, питание, социальная лента, поддержка, аватар, профиль) без полной перезагрузки страницы.

Связь клиента и сервера организована через REST API по протоколу HTTP(S) в формате JSON. Модуль auth.js выполняет регистрацию и вход, сохраняет токен и сведения о пользователе, при успешной аутентификации подгружает данные с сервера и синхронизирует ключи localStorage с таблицей user\_storage. Предметные модули (app.js, nutrition.module.js, social.module.js и др.) формируют пользовательский интерфейс, сохраняют рабочее состояние локально и при необходимости обращаются к защищённым маршрутам API с JWT в заголовке Authorization. Таким образом, клиент обеспечивает отзывчивость интерфейса и работу без постоянных запросов к серверу для каждого действия, а сервер -- единое хранилище учётных записей и межустройственную синхронизацию данных.

Клиентская часть имеет модульную структуру: авторизация отделена от тренировочной логики, а разделы питания, поддержки, аватара, социальной ленты и профиля вынесены в файлы *.module.js. Планы тренировок, история подходов, записи дневника питания и параметры интерфейса хранятся в localStorage и при наличии сети синхронизируются с сервером. Для аналитики нагрузки и прогресса используется библиотека Chart.js; справочные материалы раздела поддержки подгружаются из локального JSON-файла mental-advice.json.

\textbf{Бизнес-обоснование проекта.} Целевая аудитория -- пользователи 18--35 лет, занимающиеся в зале или дома без постоянного тренера. Ценностное предложение: единый веб-сервис <<план -- выполнение -- прогресс -- питание -- мотивация>>, тогда как аналоги закрывают отдельные задачи. MVP включает регистрацию, персональный план, журнал подходов, визуализацию нагрузки и облачную синхронизацию. Показатели эффективности: число регистраций, доля пользователей с $\geq$~3 зафиксированными тренировками за 30 дней, удержание (retention) на 7-й день. Модель монетизации (перспектива): freemium -- базовый функционал бесплатно, расширенные шаблоны программ и углублённая аналитика по подписке.

\emph{Цель настоящей работы} -- разработка и внедрение клиентской подсистемы веб-приложения для индивидуального планирования тренировок и отслеживания прогресса, обеспечивающей удобный пользовательский интерфейс, модульную организацию клиентского кода и взаимодействие с REST API серверной части. Для достижения цели необходимо решить \emph{следующие задачи:}
\begin{itemize}
\item провести анализ предметной области и сопоставить архитектурные подходы к построению клиентских веб-приложений в сфере фитнеса;
\item сформировать техническое задание и спроектировать структуру интерфейса и модулей клиентской части;
\item реализовать экраны аутентификации и профиля, клиентские модули тренировок, питания, поддержки, аватара и социальной ленты;
\item организовать локальное хранение состояния, синхронизацию с серверным хранилищем и визуализацию прогресса;
\item провести системное тестирование ключевых пользовательских сценариев в браузере.
\end{itemize}

\emph{Структура и объём работы.} Работа состоит из введения, четырёх разделов основной части, заключения, списка использованных источников, приложений и графического материала на отдельных листах. Объём текста выпускной квалификационной работы составляет \formbytotal{lastpage}{страниц}{у}{ы}{}.

\emph{Во введении} обоснована актуальность темы, описана связь клиентской и серверной подсистем, сформулированы цель, задачи и структура работы.

\emph{В первом разделе} выполнен анализ предметной области: рассмотрены потребности пользователей, классы аналогов, требования к данным и безопасности.

\emph{Во втором разделе} приведено техническое задание: назначение системы, функциональные требования, роли пользователей, сценарии использования и требования к интерфейсу клиента.

\emph{В третьем разделе} представлен технический проект клиентской части: архитектура одностраничного приложения, обоснование выбора технологий (HTML, CSS, JavaScript), описание модулей, работы с DOM, localStorage и событийной моделью, сопоставление требований с компонентами интерфейса.

\emph{В четвёртом разделе} описан рабочий проект клиентской подсистемы: спецификация модулей, системное тестирование интерфейса, проверка основных сценариев работы приложения в браузере.

В заключении сформулированы основные результаты и направления дальнейшего развития.

В приложении А представлен графический материал (плакаты). В приложении Б приведены фрагменты исходного кода клиентской части (auth.js, app.js, *.module.js, фрагменты HTML и CSS).
